\documentclass[11pt]{article}

\usepackage[margin=0.8in]{geometry}
\usepackage{amsmath,amssymb}
\usepackage{enumitem}
\usepackage{xcolor}
\usepackage{tcolorbox}
\usepackage{fancyhdr}
\usepackage{lastpage}
\usepackage{hyperref}
\usepackage{microtype}

\definecolor{uwpurple}{HTML}{4B2E83}
\definecolor{uwgold}{HTML}{B7A57A}
\definecolor{softpurple}{HTML}{F4F1FA}
\definecolor{softgold}{HTML}{FBF7EA}
\definecolor{deepgold}{HTML}{8A6F2A}

\tcbuselibrary{breakable,skins}

\pagestyle{fancy}
\fancyhf{}
\lhead{\textbf{MATH 394: Probability I}}
\rhead{\textbf{Midterm Solutions}}
\cfoot{\thepage\ of \pageref{LastPage}}
\setlength{\headheight}{14pt}

\setlength{\parindent}{0pt}
\setlength{\parskip}{0.45em}

\newcommand{\course}{MATH 394: Probability I}
\newcommand{\instructor}{Arman Jahangiri}
\newcommand{\department}{Department of Mathematics}
\newcommand{\university}{University of Washington}

\newtcolorbox{problemheader}{
    colback=softpurple,
    colframe=uwpurple,
    boxrule=0.8pt,
    arc=2mm,
    left=2mm,
    right=2mm,
    top=1mm,
    bottom=1mm
}

\newtcolorbox{solutionbox}[1][] {
    breakable,
    enhanced,
    colback=softgold,
    colframe=deepgold,
    coltitle=white,
    colbacktitle=uwpurple,
    fonttitle=\bfseries,
    title={Solution #1},
    boxrule=0.85pt,
    arc=2mm,
    left=2.5mm,
    right=2.5mm,
    top=1.5mm,
    bottom=1.5mm,
    attach boxed title to top left={xshift=2.5mm,yshift=-2mm},
    boxed title style={
        colback=uwpurple,
        colframe=uwpurple,
        arc=1.5mm,
        boxrule=0pt,
        left=1.6mm,
        right=1.6mm,
        top=0.6mm,
        bottom=0.6mm
    }
}

\newtcolorbox{finalanswerbox}[1][] {
    colback=white,
    colframe=uwpurple,
    boxrule=0.75pt,
    arc=2mm,
    left=2mm,
    right=2mm,
    top=1mm,
    bottom=1mm,
    title={\textbf{Final Answer #1}},
    coltitle=uwpurple,
    colbacktitle=softpurple,
    fonttitle=\bfseries
}

\begin{document}

\begin{titlepage}
\centering
\vspace*{1.2cm}

{\Large \textbf{\university}}\\
{\large \department}

\vspace{1.4cm}

{\Huge \color{uwpurple}\textbf{Midterm Examination Solutions}}\\[0.25cm]
{\Large \textbf{\course}}\\[0.25cm]
{\large Summer 2026 \quad | \quad Instructor: \instructor}

\vspace{1cm}

\begin{tcolorbox}[
    width=0.78\textwidth,
    colback=softpurple,
    colframe=uwpurple,
    boxrule=0.9pt,
    arc=3mm
]
\centering
\textbf{Exam date:} July 22, 2026

\textbf{Answering time:} 60 minutes

\textbf{Total:} 60 points
\end{tcolorbox}

 
 
\end{titlepage}

% ------------------------------------------------------------
\begin{problemheader}
\textbf{Problem 1. } \hfill \textbf{5 points}
\end{problemheader}
A student organization has 7 mathematics graduate students, 5 statistics graduate students, and 6 undergraduate students. A committee of 5 students is selected uniformly at random from all committees of 5.

Find the probability that the committee contains exactly 2 mathematics graduate students, exactly 1 statistics graduate student, and exactly 2 undergraduate students.

Give an exact answer. A correct expression involving binomial coefficients is acceptable.
\begin{solutionbox}
There are
\[
\binom{18}{5}
\]
possible committees of 5 students, and all are equally likely.

To form a committee of the required type, choose:
\begin{itemize}[leftmargin=*]
    \item 2 of the 7 mathematics graduate students,
    \item 1 of the 5 statistics graduate students, and
    \item 2 of the 6 undergraduate students.
\end{itemize}
Thus the number of favorable committees is
\[
\binom{7}{2}\binom{5}{1}\binom{6}{2}.
\]
Therefore,
\[
P(\text{required composition})
=
\frac{\binom{7}{2}\binom{5}{1}\binom{6}{2}}{\binom{18}{5}}.
\]
Numerically,
\[
\binom{7}{2}\binom{5}{1}\binom{6}{2}=21\cdot 5\cdot 15=1575,
\qquad
\binom{18}{5}=8568,
\]
so
\[
\frac{1575}{8568}=\frac{525}{2856}=\frac{175}{952}.
\]

\end{solutionbox}

\newpage

% ------------------------------------------------------------
\begin{problemheader}
\textbf{Problem 2.} \hfill \textbf{10 points}
\end{problemheader}

A spam filter is designed by looking at commonly occurring phrases in spam. Suppose that $80\%$ of email is spam. In $10\%$ of the spam emails, the phrase ``free money'' is used, whereas this phrase is only used in $1\%$ of non-spam emails. A new email has just arrived, which does mention ``free money.'' What is the probability that it is spam?

\begin{solutionbox}
\textbf{(a)} Let $S$ be the event that the email is spam, and let $F$ be the event that the email contains the phrase ``free money.'' We are given
\[
P(S)=0.80,\qquad P(F\mid S)=0.10,\qquad P(F\mid S^c)=0.01.
\]
Bayes' rule gives
\[
P(S\mid F)
=
\frac{P(F\mid S)P(S)}{P(F\mid S)P(S)+P(F\mid S^c)P(S^c)}.
\]
Substituting the given values,
\[
P(S\mid F)
=
\frac{0.10\cdot 0.80}{0.10\cdot 0.80+0.01\cdot 0.20}
=
\frac{0.08}{0.082}
=
\frac{40}{41}.
\]



\end{solutionbox}

\newpage
% ------------------------------------------------------------
\begin{problemheader}
\textbf{Problem 3. } \hfill \textbf{10 points}
\end{problemheader}

Let $A$ and $B$ be events.

\begin{enumerate}[label=(\alph*),leftmargin=*]
    \item Let $S$ be a sample space and $A_1,\dots, A_n\in  \mathcal{F}$ be events. Explain what it means for $A_i's$ to for a partition of $S$,  \hfill \textbf{(3 points)}
    \item Suppose $P(A) = 0.6$ and $P(B)=0.5$. Can $A$ and $B$ be disjoint? \hfill \textbf{(3 points)}
   
    \item Suppose that $A$ is independent of itself. Use the definition of independence to conclude that $P(A)=0$ or $P(A)=1$. \hfill \textbf{(4 points)}
\end{enumerate}

\begin{solutionbox}

\textbf{(a)} The events $A_1,\ldots,A_n$ form a partition of the sample space $S$ if they are pairwise disjoint,
\[
A_i\cap A_j=\emptyset,\qquad i\neq j,
\]
and
\[
\bigcup_{i=1}^n A_i=S.
\]

\textbf{(b)} If $A$ and $B$ were disjoint, then
\[
1\geq P(A\cup B)=P(A)+P(B)=0.6+0.5=1.1,
\]
which is impossible since probabilities cannot exceed $1$. Hence $A$ and $B$ cannot be disjoint.

\textbf{(c)} If $A$ is independent of itself, then
\[
P(A\cap A)=P(A)P(A).
\]
Since $A\cap A=A$,
\[
P(A)=P(A)^2.
\]
Let $p=P(A)$. Then
\[
p(1-p)=0,
\]
so $p=0$ or $p=1$.

\begin{finalanswerbox}
(a) $\{A_i\}_{i=1}^n$ is a partition iff the events are pairwise disjoint and their union is $S$.

(b) No.

(c)
\[
 {P(A)=0\quad\text{or}\quad P(A)=1.}
\]
\end{finalanswerbox}
\end{solutionbox}
\newpage
% ------------------------------------------------------------
\begin{problemheader}
\textbf{Problem 4. } \hfill \textbf{15 points}
\end{problemheader}

Suppose that $X$ has probability density function
\[
f_X(x)=
\begin{cases}
 c\,x(2-x), & 0<x<2,\\
 0, & \text{otherwise},
\end{cases}
\]
where $c$ is a constant.

\begin{enumerate}[label=(\alph*),leftmargin=*,itemsep=0.7em]
    \item Find $c$. \hfill \textbf{(3 points)}
    \item Find the cumulative distribution function $F_X(x)$ for all real $x$. \hfill \textbf{(5 points)}
    \item Find $P(X>3/2)$. \hfill \textbf{(2 points)}
    \item Find $\mathbb{E}[X]$ and $\operatorname{Var}(X)$. \hfill \textbf{(5 points)}
\end{enumerate}

\begin{solutionbox}
\textbf{(a)} A density must integrate to $1$, so
\begin{align*}
1
&=c\int_0^2 x(2-x)\,dx\\
&=c\int_0^2(2x-x^2)\,dx\\
&=c\left[x^2-\frac{x^3}{3}\right]_0^2
=c\left(4-\frac83\right)
=c\frac43.
\end{align*}
Hence
\[
c=\frac34.
\]

\textbf{(b)} For $0<x<2$,
\begin{align*}
F_X(x)
&=\int_0^x \frac34t(2-t)\,dt\\
&=\frac34\left[t^2-\frac{t^3}{3}\right]_0^x\\
&=\frac{3x^2}{4}-\frac{x^3}{4}.
\end{align*}
Therefore,
\[
F_X(x)=
\begin{cases}
0, & x\leq 0,\\[2mm]
\displaystyle \frac{3x^2-x^3}{4}, & 0<x<2,\\[2mm]
1, & x\geq 2.
\end{cases}
\]

\textbf{(c)} Using the CDF,
\begin{align*}
P\left(X>\frac32\right)
&=1-F_X\left(\frac32\right)\\
&=1-\frac{3(3/2)^2-(3/2)^3}{4}\\
&=1-\frac{27}{32}
=\frac5{32}.
\end{align*}

\textbf{(d)} First,
\begin{align*}
\mathbb{E}[X]
&=\int_0^2 x\frac34x(2-x)\,dx\\
&=\frac34\int_0^2(2x^2-x^3)\,dx\\
&=\frac34\left[\frac{2x^3}{3}-\frac{x^4}{4}\right]_0^2
=1.
\end{align*}
Also,
\begin{align*}
\mathbb{E}[X^2]
&=\int_0^2 x^2\frac34x(2-x)\,dx\\
&=\frac34\int_0^2(2x^3-x^4)\,dx\\
&=\frac34\left[\frac{x^4}{2}-\frac{x^5}{5}\right]_0^2\\
&=\frac65.
\end{align*}
Thus
\[
\operatorname{Var}(X)=\mathbb{E}[X^2]-\bigl(\mathbb{E}[X]\bigr)^2
=\frac65-1=\frac15.
\]

\begin{finalanswerbox}
\[
 {c=\frac34},\qquad
 {P(X>3/2)=\frac5{32}},\qquad
 {\mathbb{E}[X]=1},\qquad
 {\operatorname{Var}(X)=\frac15}.
\]
The CDF is
\[
 {
F_X(x)=
\begin{cases}
0, & x\leq 0,\\
(3x^2-x^3)/4, & 0<x<2,\\
1, & x\geq 2.
\end{cases}}
\]
\end{finalanswerbox}
\end{solutionbox}
 

\newpage

% ------------------------------------------------------------
\begin{problemheader}
\textbf{Problem 5.} \hfill \textbf{15 points}
\end{problemheader}

A company has found that \(40\%\) of customers who visit its website make a purchase. Suppose that \(200\) customers visit the website independently on a particular day.

Let \(X\) denote the number of customers who make a purchase.

\begin{enumerate}[label=\arabic*.,leftmargin=*,itemsep=1.2em]

\item Determine whether a \textbf{Poisson approximation} or a
\textbf{normal approximation} is more appropriate by verifying the
required conditions. Then clearly state the approximating distribution;
that is, provide \(\lambda\) if a Poisson approximation is appropriate,
or provide \(\mu\) and \(\sigma^2\) if a normal approximation is
appropriate.
\hfill \textbf{(5 points)}

\item Use the approximation from part (1) to estimate each of the following:
\hfill \textbf{(10 points)}

\begin{enumerate}[label=(\alph*),leftmargin=*,itemsep=0.7em]

    \item
    \[
    \mathbb{P}(70<X<90)
    \]
    \hfill \textbf{(3 points)}

    \item
    \[
    \mathbb{P}(X>90)
    \]
    \hfill \textbf{(3 points)}

    \item
    \[
    \mathbb{P}(X=80)
    \]
    \hfill \textbf{(4 points)}

\end{enumerate}

\end{enumerate}


\begin{solutionbox}

Since each customer independently makes a purchase with probability
\(p=0.40\), we have
\[
X\sim\operatorname{Bin}(200,0.40).
\]

\textbf{(1)} For a Poisson approximation to be appropriate, \(n\) should be large and
\(p\) should be small. Although
\[
n=200
\]
is large, \(p=0.40\) is not small. Therefore, a Poisson approximation is
not appropriate.

For a normal approximation, we verify that
\[
np=200(0.40)=80
\]
and
\[
n(1-p)=200(0.60)=120.
\]

Since both quantities are greater than \(10\), the normal approximation is
appropriate.

The mean and variance of \(X\) are
\[
\mu=np=80
\]
and
\[
\sigma^2=np(1-p)=200(0.40)(0.60)=48.
\]

Therefore, we approximate \(X\) by
\[
Y\sim N(80,48),
\]
where
\[
\sigma=\sqrt{48}\approx6.93.
\]


\textbf{(2)}

\textbf{(a)} Using the continuity correction,
\[
\mathbb{P}(70<X<90)
\approx
\mathbb{P}(70.5<Y<89.5).
\]

Standardizing,
\[
\begin{aligned}
\mathbb{P}(70.5<Y<89.5)
&=
\mathbb{P}\left(
\frac{70.5-80}{\sqrt{48}}
<
Z
<
\frac{89.5-80}{\sqrt{48}}
\right)\\
&\approx
\mathbb{P}(-1.37<Z<1.37).
\end{aligned}
\]

Using the standard normal table,
\[
\Phi(1.37)=0.9147
\]
and
\[
\Phi(-1.37)=1-\Phi(1.37)=0.0853.
\]

Therefore,
\[
\mathbb{P}(70<X<90)
\approx
0.9147-0.0853
=
0.8294.
\]


\textbf{(b)} Using the continuity correction,
\[
\mathbb{P}(X>90)
\approx
\mathbb{P}(Y>90.5).
\]

Standardizing,
\[
\begin{aligned}
\mathbb{P}(Y>90.5)
&=
\mathbb{P}\left(
Z>
\frac{90.5-80}{\sqrt{48}}
\right)\\
&\approx
\mathbb{P}(Z>1.52).
\end{aligned}
\]

Using the standard normal table,
\[
\Phi(1.52)=0.9357.
\]

Thus,
\[
\mathbb{P}(X>90)
\approx
1-\Phi(1.52)
=
1-0.9357
=
0.0643.
\]


\textbf{(c)} Using the continuity correction,
\[
\mathbb{P}(X=80)
\approx
\mathbb{P}(79.5<Y<80.5).
\]

Standardizing,
\[
\begin{aligned}
\mathbb{P}(79.5<Y<80.5)
&=
\mathbb{P}\left(
\frac{79.5-80}{\sqrt{48}}
<
Z
<
\frac{80.5-80}{\sqrt{48}}
\right)\\
&\approx
\mathbb{P}(-0.07<Z<0.07).
\end{aligned}
\]

Using the standard normal table,
\[
\Phi(0.07)=0.5279
\]
and
\[
\Phi(-0.07)=0.4721.
\]

Therefore,
\[
\mathbb{P}(X=80)
\approx
0.5279-0.4721
=
0.0558.
\]

\end{solutionbox}
 
 




\end{document}
